\section*{Lời nói đầu}
\addcontentsline{toc}{section}{Lời nói đầu}

Các em học sinh thân mến,

\medskip
Tích phân là một trong những công cụ mạnh mẽ nhất của Giải tích, nhưng ý nghĩa thực sự của nó chỉ được bộc lộ trọn vẹn khi ta nhìn nó dưới góc độ \textbf{hình học}: diện tích của một miền phẳng, thể tích của một khối không gian. Tài liệu này được biên soạn nhằm giúp các em:

\begin{itemize}[leftmargin=1.2cm]
  \item Hiểu bản chất hình học của tích phân xác định, thay vì chỉ ghi nhớ công thức một cách máy móc;
  \item Thành thạo các kỹ thuật xử lý dấu giá trị tuyệt đối, thiết lập cận tích phân từ đồ thị hoặc từ bài toán thực tế;
  \item Làm quen với ba định dạng câu hỏi mới trong kỳ thi tốt nghiệp THPT: trắc nghiệm nhiều lựa chọn, trắc nghiệm Đúng/Sai và trắc nghiệm trả lời ngắn.
\end{itemize}

\begin{hopchuy}[Hướng dẫn sử dụng tài liệu]
Mỗi bài học được tổ chức theo 5 \textbf{Dạng toán} tăng dần độ khó -- từ bài toán có sẵn dữ kiện, đến bài toán đọc đồ thị, và cuối cùng là bài toán mô hình hóa thực tiễn. Các em nên:
\begin{enumerate}[leftmargin=1.2cm]
  \item Đọc kỹ hộp \textcolor{teal}{\textbf{Định lý}} trước khi làm ví dụ;
  \item Tự trình bày lời giải trước khi xem \textcolor{navy}{\textbf{Lời giải}} mẫu;
  \item Ghi chú các \textcolor{amber}{\textbf{Chú ý}} về sai lầm thường gặp -- đây là những "bẫy" hay xuất hiện trong đề thi.
\end{enumerate}
\end{hopchuy}

\bigskip
\hfill \textit{Thầy Tâm dạy Toán}

\newpage
